\documentclass[12pt]{article}
\usepackage[spanish]{babel}
\usepackage[utf8]{inputenc}
\usepackage[T1]{fontenc}
\usepackage{amsmath, amssymb}
\usepackage{geometry}
\usepackage{enumitem}
\usepackage{needspace}
\usepackage{tikz}
\usetikzlibrary{babel}

\geometry{letterpaper, margin=2cm}
\setlength{\parindent}{0pt}
\setlength{\parskip}{6pt}

\newcommand{\puntografico}{0.09}

\begin{document}

\begin{center}
  {\Large \textbf{Borrador de items}}\\
  {\large Bloque 1. Geometria - Afirmacion 5}\\[4pt]
  Prueba Nacional Estandarizada de Matematica, Secundaria 2026
\end{center}

\section*{Afirmacion}
Resuelve problemas relacionados con secciones planas o elementos de figuras geometricas tridimensionales.

\section*{Items propuestos}

\begin{enumerate}[label=\textbf{\arabic*.}, itemsep=1.05cm]

\Needspace{0.42\textheight}
\item
En un taller se elaboran recuerdos a partir de esferas de acrilico. A cada esfera se le realiza un corte plano que pasa por el centro de la esfera.

De acuerdo con la informacion anterior, la seccion plana obtenida en cada recuerdo, producto del corte realizado a esa esfera, corresponde a una

\begin{enumerate}[label=\Alph*)]
  \item circunferencia.
  \item hiperbola.
  \item parabola.
  \item recta.
\end{enumerate}

\Needspace{0.48\textheight}
\item
Considere la siguiente informacion:

Una persona artesana trabaja con un bloque de cera con forma de cilindro circular recto, cuyo radio de la base mide $8$ cm y la altura del cilindro mide $24$ cm. Para elaborar una pieza decorativa, realiza un corte plano al bloque, que pasa por el centro de las bases y es perpendicular a estas.

De acuerdo con la informacion anterior, \textquestiondown cual es el perimetro de la seccion plana obtenida producto del corte realizado a ese bloque?

\begin{enumerate}[label=\Alph*)]
  \item $32$ cm
  \item $64$ cm
  \item $80$ cm
  \item $192$ cm
\end{enumerate}

\Needspace{0.78\textheight}
\item
Maria debe hacer cinco sombreros de cumpleanos como proyecto para un curso de fiestas infantiles. Ella confecciono los sombreros con un diametro de $14$ cm y una altura de $25$ cm. Si quiere pegar cinta formando una circunferencia paralela a la base y a $10$ cm de altura sobre la base del cono, como se muestra en la siguiente figura, \textquestiondown cuanta cinta debe comprar, aproximadamente, para los cinco sombreros?

\begin{center}
\begin{tikzpicture}[scale=0.62]
  \path[use as bounding box] (-4.4,-0.9) rectangle (4.7,7.3);
  \coordinate (V) at (0,6);
  \coordinate (L) at (-3,0);
  \coordinate (R) at (3,0);
  \coordinate (C) at (0,0);
  \coordinate (S) at (0,2.4);
  \coordinate (P) at (1.8,2.4);
  \draw[very thick] (V) -- (L);
  \draw[very thick] (V) -- (R);
  \draw[very thick] (0,0) ellipse (3 and 0.55);
  \draw[gray!40, very thick] (0,2.4) ellipse (1.8 and 0.33);
  \draw[dashed, thick] (0,0) -- (0,6);
\end{tikzpicture}
\end{center}

\begin{enumerate}[label=\Alph*)]
  \item $42$ cm
  \item $84$ cm
  \item $132$ cm
  \item $220$ cm
\end{enumerate}

\end{enumerate}

\section*{Clave y justificacion breve}

\begin{enumerate}[label=\textbf{\arabic*.}]
  \item \textbf{A.} Si un plano corta una esfera pasando por su centro, la seccion plana obtenida corresponde a una circunferencia.
  \item \textbf{C.} Un corte perpendicular a las bases y que pasa por sus centros produce un rectangulo cuyos lados miden la altura del cilindro, $24$ cm, y el diametro de la base, $16$ cm. Por eso, el perimetro es $2(24+16)=80$ cm.
  \item \textbf{C.} A $10$ cm sobre la base quedan $15$ cm hasta el vertice, de un total de $25$ cm. Por semejanza, el radio de la circunferencia de la cinta es $7\cdot\frac{15}{25}=4,2$ cm. La longitud de una cinta es aproximadamente $2\pi(4,2)\approx 26,4$ cm; para cinco sombreros se necesitan aproximadamente $132$ cm.
\end{enumerate}

\end{document}
